\documentclass[11pt]{article}

\usepackage[margin=1in]{geometry}
\usepackage[T1]{fontenc}
\usepackage[utf8]{inputenc}
\usepackage{lmodern}
\usepackage{microtype}
\usepackage{amsmath, amssymb, amsthm}
\usepackage{graphicx}
\usepackage{booktabs}
\usepackage{hyperref}
\usepackage{xcolor}
\usepackage{enumitem}

\hypersetup{
  colorlinks=true,
  linkcolor=blue!60!black,
  citecolor=blue!60!black,
  urlcolor=blue!60!black
}

\title{Inverse Planning in a Simple Gridworld\\
\large Exact Bayesian Inference and RNN-Based Approximations}
\author{CS6208 Project Report Draft}
\date{\today}

\begin{document}

\maketitle

\begin{abstract}
This report studies inverse planning in a simple gridworld where an observer
infers an agent's latent goal from its actions. We use an exact
approximately-rational planning model as a reference baseline, then evaluate
recurrent neural network variants that either predict goals directly or
approximate inverse planning through a goal-conditioned action model.
\end{abstract}

\section{Introduction}
Inverse planning frames intention understanding as Bayesian inference over
latent mental states \cite{baker2009action, baker2017rational}. Given a model
of how an approximately rational agent acts toward a goal, an observer can
invert that model to infer which goal best explains an observed trajectory.

This project focuses on a simple fully observed gridworld based on the
\texttt{inv\_plan\_from\_scratch} task. The main reason for choosing this
setting is methodological: exact inverse planning remains tractable, allowing us
to construct a clean baseline and generate reliable supervision for learned
models. On top of this baseline, we study several RNN-based approximations that
may be useful when exact inference becomes expensive.

\paragraph{Project contributions.}
The report is structured around three contributions:
\begin{enumerate}[leftmargin=1.5em]
  \item an exact inverse-planning baseline for a simple gridworld;
  \item a reusable data-generation pipeline that labels trajectories with exact
  posterior beliefs over goals;
  \item a comparison of direct goal-prediction and goal-conditioned policy RNN
  variants.
\end{enumerate}

\section{Related Work}
Bayesian inverse planning and related theory-of-mind models have been used to
infer goals, beliefs, and other mental states from observed behavior
\cite{baker2009action, baker2017rational}. Probabilistic programming systems
such as Gen support flexible implementations of such models while preserving a
clean separation between model specification and inference \cite{cusumano2018gen}.

For more structured planning settings, partially observable models and
belief-space inference introduce substantial computational challenges
\cite{kaelbling1998planning}. Recent code resources such as \texttt{memo} and
LaBToM provide useful implementation references for richer recursive or
belief-based planning settings. On the function-approximation side, recurrent
neural networks and convolutional recurrent models offer a natural way to map
trajectory observations to latent-state predictions \cite{hochreiter1997lstm,
cho2014gru}. Our project uses these learned models as approximations to exact
Bayesian inference rather than as replacements for the underlying behavioral
theory.

\section{Problem Setup}
\subsection{Gridworld Environment}
The environment is a small deterministic gridworld with walls, a fixed initial
location, and a discrete set of candidate goals. At each time step, the agent
chooses among five actions: stay, up, down, left, and right. Invalid moves
leave the agent in place.

\subsection{Approximately Rational Agent Model}
We assume that the agent acts approximately rationally with respect to a latent
goal $g$. Let $Q_g(s, a)$ denote the state-action value of taking action $a$ in
state $s$ under goal $g$. The agent follows a Boltzmann policy
\begin{equation}
P(a_t \mid s_t, g) =
\frac{\exp(\beta Q_g(s_t, a_t))}
{\sum_{a' \in \mathcal{A}} \exp(\beta Q_g(s_t, a'))},
\end{equation}
where $\beta$ controls how strongly behavior concentrates on high-value
actions.

\subsection{Inverse Planning Objective}
Given an observed action sequence $a_{1:T}$, the observer seeks the posterior
distribution over goals:
\begin{equation}
P(g \mid a_{1:T}) \propto P(g) \prod_{t=1}^{T} P(a_t \mid s_t, g).
\end{equation}
Because the environment is deterministic and the goal set is small, this
posterior can be computed exactly by enumerating over goals.

\section{Exact Inverse-Planning Baseline}
\subsection{Planning via Shortest Paths}
In the simple gridworld, action values can be induced from shortest-path
distances to the goal. This makes planning transparent and inexpensive compared
with general POMDP belief-space planning.

\subsection{Exact Posterior Inference}
For each candidate goal, we evaluate the likelihood of the observed action
sequence under the corresponding Boltzmann policy and normalize across goals.
This yields an exact posterior over goals at the final time step and also
supports online posterior updates over time.

\subsection{Role as a Reference Model}
The exact model serves two purposes: it provides a principled baseline for
evaluation, and it supplies high-quality labels for supervised learning. This
is important because learned models can then be assessed not only on final-goal
accuracy but also on how well they match the posterior dynamics of the exact
observer.

\section{Dataset Construction}
\subsection{Trajectory Generation}
We generate trajectories by sampling a latent goal and then rolling out the
approximately rational policy for a fixed horizon. Each trajectory includes the
visited states, chosen actions, and action probabilities along the rollout.

\subsection{Observation Encoding}
Each time step is encoded as a spatial tensor containing the wall map, current
agent position, goal markers, and a time feature. This representation is chosen
to support both simple recurrent baselines and convolutional recurrent models.

\subsection{Supervision Targets}
For each trajectory, we store:
\begin{itemize}[leftmargin=1.5em]
  \item the true latent goal;
  \item the observed action sequence;
  \item the exact online posterior over goals at every step.
\end{itemize}
These targets support both direct goal prediction and imitation-style action
modeling.

\section{RNN-Based Approximations}
\subsection{Action-History GRU}
The first baseline encodes observations and previous actions into a GRU and
predicts the goal directly. This model provides a simple learned approximation
with relatively weak spatial inductive bias.

\subsection{Convolutional GRU Classifier}
The second baseline first applies a convolutional encoder to each gridworld
frame before the recurrent update. The motivation is to better preserve the
spatial structure and local invariances of the environment.

\subsection{Goal-Conditioned Policy RNN}
The third model conditions on a candidate goal and predicts actions over time.
At test time, the observer can score the observed trajectory under each
candidate goal and normalize these scores into a posterior distribution. This
approach is closer in spirit to inverse planning because it models action
generation rather than predicting the goal in one step.

\section{Experimental Setup}
\subsection{Training Objectives}
Direct classifiers are trained to predict the true goal and may optionally be
regularized to match the exact online posterior. The goal-conditioned policy is
trained with an action-prediction loss under the true goal.

\subsection{Evaluation Metrics}
We plan to report final-goal accuracy, posterior divergence from the exact
observer, and qualitative trajectory-level comparisons. These metrics together
measure both discrete prediction quality and the fidelity of online inference.

\subsection{Implementation Notes}
The codebase is organized so the exact inverse-planning pipeline, data
generation, and learned models are reusable from both scripts and the demo
notebook. This separation makes it straightforward to replace the exact
handcrafted planner with a \texttt{memo}-based backend in future work.

\section{Results}
This section will present the final quantitative and qualitative results.

\subsection{Exact Baseline Behavior}
Planned material:
\begin{itemize}[leftmargin=1.5em]
  \item example trajectories under different goals;
  \item posterior updates over time for representative sequences;
  \item discussion of the effect of $\beta$ on inference sharpness.
\end{itemize}

\subsection{Comparison with RNN Variants}
Planned material:
\begin{itemize}[leftmargin=1.5em]
  \item a summary table of final-goal accuracy;
  \item posterior-tracking comparisons against the exact baseline;
  \item failure cases where learned models diverge from exact inference.
\end{itemize}

\section{Discussion}
The main advantage of the current project design is that it is methodologically
clean: the exact model remains simple enough to understand and verify, while
the learning problem is still nontrivial. The main limitation is that the
project does not yet solve Bayesian inverse planning in belief space for a true
POMDP, which is considerably harder.

This staged design is still useful. It establishes a solid baseline story and a
reusable code path for richer future models. In particular, a next step is to
swap the handcrafted planner for a \texttt{memo}-based or other probabilistic
planning backend that can better support partial observability and recursive
reasoning.

\section{Conclusion}
This project studies inverse planning in a simple gridworld using an exact
approximately-rational observer and several RNN-based approximations. The exact
model provides both a strong baseline and a data source, while the learned
models test how far direct sequence modeling can approximate Bayesian goal
inference in a controlled setting.

\bibliographystyle{plain}
\bibliography{refs}

\end{document}
